% Prozentrechnung – bank/prozentrechnung/ (zusammenbau v0.4, Heft)
\einheitenkopf[t3]{Prozentrechnung}
\setcounter{aufgabe}{22}

% Anlauf (ohne Prüfkennung)
\begin{aufgabe}{Umwandeln – Anlauf}
\begin{teile}
\teil Wie viel Prozent sind grau?

\streifenfeld[2]{50}{$0\,\%$}{$100\,\%$}
\teil $\frac{37}{100}$ – wie viel Prozent? \leerfeld[\%]
\end{teile}
\end{aufgabe}

\begin{aufgabe}{Umwandeln – Prüfungsaufgaben}
\begin{teile}
\teil $5\,\%$ der Pakete kommen zu spät – welche Aussage passt? Kreuze an. \hfill (P10 2020 OS)\\ \kreuz{$5$ von $10$ Paketen kommen zu spät.}\\ \kreuz{Jedes 5. Paket kommt zu spät.}\\ \kreuz{$5$ von $100$ Paketen kommen zu spät.}
\teil $8\,\%$ der Gäste bestellen Tee – welche Aussage passt? Kreuze an. \hfill (P10 2020 OS)\\ \kreuz{Jeder 8. Gast bestellt Tee.}\\ \kreuz{$8$ von $100$ Gästen bestellen Tee.}\\ \kreuz{$8$ von $10$ Gästen bestellen Tee.}
\teil Jeder vierte Besucher kauft ein Programmheft – wie viel Prozent? Kreuze an. \hfill (P10 2022 OS)\\ \kreuz{$4\,\%$}\\ \kreuz{$25\,\%$}\\ \kreuz{$40\,\%$}\\ \kreuz{$75\,\%$}
\teil Jedes zwanzigste Ei einer Lieferung ist zerbrochen – wie viel Prozent? Kreuze an. \hfill (P10 2022 OS)\\ \kreuz{$2\,\%$}\\ \kreuz{$5\,\%$}\\ \kreuz{$20\,\%$}\\ \kreuz{$80\,\%$}
\end{teile}
\end{aufgabe}

\begin{aufgabe}{Umwandeln – Prüfungsaufgaben – weiter}
\begin{teile}
\teil Schulweg der Kinder einer Schule. Genau eine Aussage ist falsch. Kreuze sie an und schreibe sie richtig. \hfill (P10 2023 OS)\\ \kreuz{Aussage 1: Ein Viertel der Kinder fährt mit dem Bus.}\\ \kreuz{Aussage 2: Die Hälfte der Kinder fährt mit dem Rad.}

\sachtabelle{lr}{Schulweg & Kinder}{Bus & 150\\ Rad & 200\\ zu Fuß & 180\\ Auto & 70\\ gesamt & 600}
\teil Mitglieder eines Sportvereins. Genau eine Aussage ist falsch. Kreuze sie an und schreibe sie richtig. \hfill (P10 2023 OS)\\ \kreuz{Aussage 1: Jedes fünfte Mitglied ist ein Kind.}\\ \kreuz{Aussage 2: Die Hälfte der Mitglieder sind Erwachsene.}

\sachtabelle{lr}{Gruppe & Mitglieder}{Kinder & 600\\ Jugendliche & 480\\ Erwachsene & 1\,200\\ Senioren & 120\\ gesamt & 2\,400}
\teil Freizeit: Anteil der Befragten in der Stadt und auf dem Land. Genau eine Aussage ist falsch. Kreuze sie an und schreibe sie richtig. \hfill (P10 2019 OS)\\ \kreuz{Aussage 1: Ein Viertel der Befragten auf dem Land geht in die Bibliothek.}\\ \kreuz{Aussage 2: Mehr als die Hälfte der Befragten auf dem Land treibt Sport im Verein.}

\sachtabelle{lcc}{Freizeit & Stadt & Land}{Sportverein & 40\,\% & 55\,\%\\ Kino & 12\,\% & 5\,\%\\ Bibliothek & 25\,\% & 10\,\%\\ Schwimmbad & 60\,\% & 48\,\%}
\teil Frühstück zu Hause: Anteil der Kinder an der Grundschule und an der Oberschule. Genau eine Aussage ist falsch. Kreuze sie an und schreibe sie richtig. \hfill (P10 2019 OS)\\ \kreuz{Aussage 1: Drei von fünf Oberschülern frühstücken immer zu Hause.}\\ \kreuz{Aussage 2: Jeder zehnte Grundschüler frühstückt nie zu Hause.}

\sachtabelle{lcc}{Frühstück & Grundschule & Oberschule}{immer & 80\,\% & 60\,\%\\ manchmal & 15\,\% & 30\,\%\\ nie & 5\,\% & 10\,\%}
\steil Wie oft Jugendliche Sport treiben (Anteile). Jemand behauptet: „Mehr als die Hälfte der Jugendlichen treibt mindestens zweimal pro Woche Sport.“ Stimmt das? Begründe mit den Werten der Tabelle. \hfill (P10 2017 OS)

\sachtabelle{lr}{Sport & Anteil}{nie & 18{,}5\,\%\\ einmal pro Woche & 32{,}0\,\%\\ zweimal pro Woche & 27{,}5\,\%\\ dreimal oder öfter & 22{,}0\,\%}
\steil Haushalte nach Zahl der Autos (Anteile). Jemand behauptet: „Mehr als $70\,\%$ der Haushalte haben mindestens ein Auto.“ Stimmt das? Begründe mit den Werten der Tabelle. \hfill (P10 2017 OS)

\sachtabelle{lr}{Autos & Anteil}{kein Auto & 28{,}6\,\%\\ ein Auto & 45{,}3\,\%\\ zwei Autos & 21{,}8\,\%\\ drei oder mehr & 4{,}3\,\%}
\end{teile}
\end{aufgabe}

\begin{aufgabe}{Umwandeln – Prüfungsaufgaben – weiter}
\begin{teile}
\teil Strom einer Gemeinde nach Energiequelle – trage den fehlenden Anteil ein. \hfill (P10 2021 OS)

\sachtabelle{lr}{Energiequelle & Anteil}{Wind & 38\,\%\\ Sonne & 24\,\%\\ Biogas & \leerzelle\\ Wasser & 9\,\%\\ Sonstiges & 3\,\%\\ gesamt & 100\,\%}
\teil Fläche eines Landkreises nach Nutzung – trage den fehlenden Anteil ein. \hfill (P10 2021 OS)

\sachtabelle{lr}{Nutzung & Anteil}{Wald & 34\,\%\\ Acker & 41\,\%\\ Siedlung & \leerzelle\\ Wasser & 4\,\%\\ Sonstiges & 2\,\%\\ gesamt & 100\,\%}
\end{teile}
\end{aufgabe}

% Anlauf (ohne Prüfkennung)
\begin{aufgabe}{Prozentsatz – Anlauf}
\begin{teile}
\teil Streifen: $50$ Kinder, grau: $10$ Kinder – wie viel Prozent?

\streifenfeld{20}{$0$}{$50$ Kinder}
\teil $17$ von $100$ Schülern spielen ein Instrument – wie viel Prozent? \leerfeld[\%]
\teil $19$ von $50$ Kindern – wie viel Prozent? \leerfeld[\%]
\end{teile}
\end{aufgabe}

\begin{aufgabe}{Prozentsatz – Prüfungsaufgaben}
\begin{teile}
\teil Bei einer Tombola gibt es $21$ Nieten und $3$ Hauptgewinne. Wie viel Prozent der Lose sind Hauptgewinne? \hfill (P10 2018 OS) \\ \leerfeld[\%]
\teil Auf dem Blech liegen $11$ Muffins mit Blaubeeren und $4$ mit Kirschen. Wie viel Prozent sind Kirschmuffins? Runde auf eine Stelle. \hfill (P10 2018 OS) \\ \leerfeld[\%]
\teil Einwohner einer Stadt: gesamt $480\,000$; Kinder $72\,000$; Jugendliche $57\,000$; Erwachsene $264\,000$; Senioren $87\,000$. Wie viel Prozent sind Jugendliche? Runde auf eine Stelle. \hfill (P10 2023 OS)
\teil Wasser eines Haushalts pro Tag: gesamt $380$ l; Toilette $102$ l; Duschen $136$ l; Wäsche $46$ l; Rest $96$ l. Wie viel Prozent entfallen aufs Duschen? Runde auf eine Stelle. \hfill (P10 2023 OS)
\teil Von $64$ Bewerbungen wurden $47$ angenommen. Wie viel Prozent wurden abgelehnt? Runde auf eine Stelle. \hfill (P10 2015 OS)
\teil Ein Zug fuhr an $92$ Tagen, $79$-mal war er pünktlich. An wie viel Prozent der Tage war er verspätet? Runde auf eine Stelle. \hfill (P10 2015 OS)
\teil Eine Stadt hat $40\,000$ Einwohner, $1\,080$ sind in diesem Jahr zugezogen. Wie viel Prozent sind das? \hfill (P10 2014 OS) \\ \leerfeld[\%]
\teil Von $12\,500$ verkauften Konzertkarten wurden $175$ zurückgegeben. Wie viel Prozent sind das? \hfill (P10 2014 OS) \\ \leerfeld[\%]
\teil Körpergrößen der $13$ Spieler einer Mannschaft in cm: $172$, $185$, $169$, $190$, $178$, $181$, $166$, $188$, $175$, $183$, $170$, $192$, $177$. Wie viel Prozent sind größer als $180$ cm? Runde auf eine Stelle. \hfill (P10 2025 OS)
\teil Wartezeiten von $9$ Patienten in Minuten: $12$, $35$, $8$, $22$, $41$, $17$, $30$, $5$, $26$. Wie viel Prozent warteten länger als $20$ Minuten? Runde auf eine Stelle. \hfill (P10 2025 OS)
\end{teile}
\end{aufgabe}

% Anlauf (ohne Prüfkennung)
\begin{aufgabe}{Prozentwert – Anlauf}
\begin{teile}
\teil $50\,\%$ von $80\,€$ – wie viel?

\streifen[2]{50}{$0$}{$80\,€$}
\leerfeld[€]
\teil $6\,\%$ von $400$ kg \leerfeld[kg]
\end{teile}
\end{aufgabe}

\begin{aufgabe}{Prozentwert – Prüfungsaufgaben}
\begin{teile}
\teil Ein Fernseher kostet $480\,€$. Bei Barzahlung gibt es $15\,\%$ Rabatt. Wie viel Euro spart man? Kreuze an. \hfill (P10 2021 OS)\\ \kreuz{$32\,€$}\\ \kreuz{$408\,€$}\\ \kreuz{$72\,€$}\\ \kreuz{$552\,€$}
\teil Eine Waschmaschine kostet $650\,€$. Im Angebot gibt es $30\,\%$ Rabatt. Wie viel Euro spart man? Kreuze an. \hfill (P10 2021 OS)\\ \kreuz{$455\,€$}\\ \kreuz{$65\,€$}\\ \kreuz{$845\,€$}\\ \kreuz{$195\,€$}
\teil Ein Tablet kostet $240{,}00\,€$. Es gibt $15\,\%$ Rabatt. Wie viel Euro Rabatt sind das? \hfill (P10 2017 OS) \\ \leerfeld[€]
\teil Eine Gitarre kostet $180{,}00\,€$. Es gibt $35\,\%$ Rabatt. Wie viel Euro Rabatt sind das? \hfill (P10 2017 OS) \\ \leerfeld[€]
\teil $17\,\%$ von $60\,€$ \hfill (P10 2014 OS) \\ \leerfeld[€]
\teil $19\,\%$ von $40\,€$ \hfill (P10 2014 OS) \\ \leerfeld[€]
\teil $40\,\%$ von $85\,€$ \hfill (P10 2026 FOR) \\ \leerfeld[€]
\teil $60\,\%$ von $45\,€$ \hfill (P10 2026 FOR) \\ \leerfeld[€]
\teil Bei einer Umfrage unter $800$ Haushalten hatten $35\,\%$ ein E-Bike, sechs Jahre später $52\,\%$. Um wie viele Haushalte hat die Zahl zugenommen? \hfill (P10 2019 OS)
\teil Eine Schule hat $650$ Schüler. Im Sommer kommen $44\,\%$ mit dem Rad, im Winter $26\,\%$. Wie viele Schüler kommen im Winter weniger mit dem Rad? \hfill (P10 2019 OS)
\end{teile}
\end{aufgabe}

% Anlauf (ohne Prüfkennung)
\begin{aufgabe}{Grundwert – Anlauf}
\begin{teile}
\teil Grau sind $30\,\%$, das sind $6$ Kästchen – wie viele Kästchen hat der ganze Streifen?

\streifen[0]{30}{$0\,\%$}{$100\,\%$}
\leerfeld[Kästchen]
\teil $50\,\%$ sind $35$ kg – wie viel ist das Ganze? \leerfeld[kg]
\teil $3\,\%$ sind $12$ kg – wie viel ist das Ganze? \leerfeld[kg]
\end{teile}
\end{aufgabe}

\begin{aufgabe}{Grundwert – Prüfungsaufgaben}
\begin{teile}
\teil Beim Kauf einer Mütze spart Mira $4\,€$, das sind $25\,\%$ des alten Preises. Wie viel kostete die Mütze vorher? \hfill (P10 2025 OS) \\ \leerfeld[€]
\teil Bei einer Konzertkarte beträgt der Rabatt $15\,€$, das sind $30\,\%$ des alten Preises. Wie hoch war der alte Preis? \hfill (P10 2025 OS) \\ \leerfeld[€]
\teil Lisa bekommt $40\,\%$ einer Erbschaft, das sind $6\,000\,€$. Wie groß ist die ganze Erbschaft? \hfill (P10 2023 OS) \\ \leerfeld[€]
\teil $20\,\%$ der Ernte eines Obstbauern sind $450$ kg Äpfel. Wie groß ist die ganze Ernte? \hfill (P10 2023 OS) \\ \leerfeld[kg]
\end{teile}
\end{aufgabe}

% Anlauf (ohne Prüfkennung)
\begin{aufgabe}{Veränderung – Anlauf}
\begin{teile}
\teil um $20\,\%$ gesenkt – was bedeutet dasselbe? Kreuze an.\\ \kreuz{auf $20\,\%$ gesenkt}\\ \kreuz{auf $80\,\%$ gesenkt}\\ \kreuz{um $80\,\%$ gesenkt}
\teil $80\,€$ um $10\,\%$ erhöht – neuer Preis? \leerfeld[€]
\teil $45\,€$ um $20\,\%$ erhöht – mit Faktor \leerfeld[€]
\end{teile}
\end{aufgabe}

\begin{aufgabe}{Veränderung – Prüfungsaufgaben}
\begin{teile}
\teil Ein Liter Milch kostete $1{,}15\,€$, jetzt kostet er $1{,}29\,€$. Um wie viel Prozent ist der Preis gestiegen? Runde auf eine Stelle. \hfill (P10 2026 FOR)
\teil Eine Kugel Eis kostete $1{,}70\,€$, jetzt kostet sie $1{,}90\,€$. Um wie viel Prozent ist der Preis gestiegen? Runde auf eine Stelle. \hfill (P10 2026 FOR)
\teil Ein Kino hatte im Mai $1\,340$ Besucher, im Juni $1\,120$. Um wie viel Prozent ist die Zahl gesunken? Runde auf eine Stelle. \hfill (P10 2022 OS)
\teil Eine Bücherei verlieh im März $860$ Bücher, im April $780$. Um wie viel Prozent ist die Zahl gesunken? Runde auf eine Stelle. \hfill (P10 2022 OS)
\teil Ein Online-Shop hatte $48\,000$ Bestellungen, im Jahr darauf $57\,600$. Um wie viel Prozent ist die Zahl gestiegen? \hfill (P10 2016 OS)
\teil Eine Stadt hatte $64\,000$ Einwohner, zehn Jahre später $72\,000$. Um wie viel Prozent ist die Zahl gestiegen? \hfill (P10 2016 OS)
\teil Ein Busticket kostet $2{,}40\,€$ und wird um $25\,\%$ teurer. Wie viel kostet es dann? \hfill (P10 2024 OS) \\ \leerfeld[€]
\teil Ein Schwimmkurs kostet $65\,€$ und wird um $12\,\%$ teurer. Wie viel kostet er dann? \hfill (P10 2024 OS) \\ \leerfeld[€]
\teil Kletterhalle: Erwachsene $14\,€$, Jugendliche (bis $17$ Jahre) $9\,€$. Familie Demir: Mutter, Vater, Tochter ($15$ Jahre), Sohn ($12$ Jahre). Montags gibt es $15\,\%$ Rabatt. Wie viel zahlt die Familie am Montag? \hfill (P10 2015 OS)
\teil Theater: Erwachsene $22\,€$, Schüler $13\,€$. Es kommen Oma, Mutter, Vater und ein Schüler. Mit der Familienkarte gibt es $25\,\%$ Rabatt. Wie viel zahlen sie zusammen? \hfill (P10 2015 OS)
\teil Für Radwege gilt: Steigung höchstens $5\,\%$. Ergänze: „$5\,\%$ Steigung heißt: auf \leerfeld[m] waagerechter Strecke \leerfeld[m] Höhenunterschied.“ Herr Lenz sagt: „Ein Radweg, der auf $240$ m waagerechter Strecke um $15$ m ansteigt, ist zu steil.“ Hat er recht? Rechne. \hfill (P10 2025 OS)
\teil Eine Garagenzufahrt darf höchstens $15\,\%$ Steigung haben. Ergänze: „$15\,\%$ Steigung heißt: auf $100$ m waagerechter Strecke \leerfeld[m] Höhenunterschied.“ Eine Zufahrt überwindet auf $8$ m waagerechter Strecke $1{,}1$ m. Frau Roth sagt, sie sei zu steil. Hat sie recht? Rechne. \hfill (P10 2025 OS)
\end{teile}
\end{aufgabe}
